\documentclass[11pt]{article}
% Supplementary material for the RESS version (single-anonymized review).
% numbers from numbers.tex and ../proofs_numbers.tex (all generated from results/*.json).
\usepackage[margin=2.3cm]{geometry}
\usepackage{amsmath,amssymb,amsthm}
\usepackage{booktabs}
\usepackage{float}
\usepackage{longtable}
\usepackage[hidelinks]{hyperref}
\hypersetup{pdfauthor={Lam Quoc Dat}, pdftitle={Supplementary material: No Free Credit (RESS)}}
% generated by tools/make_numbers.py from results/*.json -- DO NOT EDIT BY HAND (checklist C6)
\newcommand{\nNzero}{59} % n0(5%,5%) = ceil(ln .05/ln .95); c2_mc.primary n0
\newcommand{\nNzeroTwenty}{14} % n0(20%,5%); exchange_rate_v2.marked.OP*.n0
\newcommand{\ncExactA}{8} % c2_mc.primary[delta'=0.075].c_integer_exact (events)
\newcommand{\ncStarA}{7} % c2_mc.primary[delta'=0.075].c_star (continuous lower bound)
\newcommand{\ncSaveA}{14\%} % c / n0 = 8/59
\newcommand{\ncExactB}{14} % c2_mc.primary[delta'=0.1].c_integer_exact (events)
\newcommand{\ncStarB}{13} % c2_mc.primary[delta'=0.1].c_star (continuous lower bound)
\newcommand{\ncSaveB}{24\%} % c / n0 = 14/59
\newcommand{\ncExactC}{22} % c2_mc.primary[delta'=0.15].c_integer_exact (events)
\newcommand{\ncStarC}{21} % c2_mc.primary[delta'=0.15].c_star (continuous lower bound)
\newcommand{\ncSaveC}{37\%} % c / n0 = 22/59
\newcommand{\ncExactD}{27} % c2_mc.primary[delta'=0.2].c_integer_exact (events)
\newcommand{\ncStarD}{27} % c2_mc.primary[delta'=0.2].c_star (continuous lower bound)
\newcommand{\ncSaveD}{46\%} % c / n0 = 27/59
\newcommand{\nCtwoNgrid}{40} % c2_mc grid points
\newcommand{\nCtwoMaxDevPP}{0.29} % c2_mc.max_abs_dev_pp_exact (percentage points, 40-point grid, MC 1e5)
\newcommand{\nCthreeA}{72} % c3_numbers q_U=0
\newcommand{\nCthreeB}{91} % q_U=1%
\newcommand{\nCthreeC}{122} % q_U=2%
\newcommand{\nPopScenesTen}{29} % c4_power.scenes_needed_c4a_pop U<=10% (J=0)
\newcommand{\nPopScenesFive}{59} % U<=5% (J=0)
\newcommand{\nLemmaMin}{0.2500} % c4_power.lemma_mean_median.min_F_floor_Np
\newcommand{\nCorpusT}{3,170} % corpus17_manifest.counts.mode_T_total (events)
\newcommand{\nCorpusTvideos}{101} % mode_T_real_videos
\newcommand{\nMainScenes}{48} % calibration_scenes.main
\newcommand{\nNightScenes}{6} % calibration_scenes.night
\newcommand{\nTurbScenes}{4} % calibration_scenes.turbulence
\newcommand{\nPTZScenes}{4} % calibration_scenes.excluded_PTZ
\newcommand{\nJitterScenes}{26} % calibration_scenes.jitter (cameraJitter, post-F1 correction v1.1)
\newcommand{\nMainCDnet}{28} % p1_background17.by_tier.main CDnet
\newcommand{\nMainLASIESTA}{20} % by_tier.main LASIESTA
\newcommand{\nMainN}{114} % e0_units main N (events, d>=32)
\newcommand{\nMainDw}{0} % e0_units main D_worst_phase.num (events)
\newcommand{\nMainDwPct}{0.00\%} % 0/114
\newcommand{\nMainDexp}{0.0000} % expected discordant events (mean over phases and 3 variants)
\newcommand{\nMainDexpPct}{0.00\%} % 0.0000/114
\newcommand{\nMainUfleet}{2.59\%} % U_CP(0,114)
\newcommand{\nMainJ}{0} % scenes with a discordant event
\newcommand{\nMainM}{48} % scenes
\newcommand{\nMainUpop}{6.05\%} % U_CP(0,48)
\newcommand{\nNightN}{25} % night events d>=32
\newcommand{\nNightDw}{12} % night worst-phase discordant
\newcommand{\nNightDwPct}{48\%} % 12/25
\newcommand{\nNightDexp}{8.21} % night expected discordant events (mean over phases and 3 variants)
\newcommand{\nNightDexpMaj}{8.31} % night expected discordant events, majority rule
\newcommand{\nNightUfleet}{65.9\%} % night U_fleet worst
\newcommand{\nNightJ}{5} % night J
\newcommand{\nNightUpop}{99.1\%} % night U_pop
\newcommand{\nJitN}{62} % jitter events d>=32
\newcommand{\nJitDw}{3} % jitter worst-phase discordant
\newcommand{\nJitUfleet}{12.0\%} % jitter U_fleet worst
\newcommand{\nJitJ}{1} % jitter J
\newcommand{\nJitM}{26} % jitter scenes
\newcommand{\nJitUpop}{17.0\%} % jitter U_pop
\newcommand{\nJitFC}{0} % jitter cameras falsely certified, main-only calibration, phase-averaged (OP*)
\newcommand{\nJitFCch}{0} % same at (16,16)
\newcommand{\nJitPgt}{0} % jitter cameras with phase-averaged p_s > eps (OP*)
\newcommand{\nTurbN}{2} % turbulence events d>=32 (scenes run: 2)
\newcommand{\nAllN}{203} % all tiers N
\newcommand{\nAllDw}{16} % all tiers worst-phase discordant
\newcommand{\nAllUfleet}{11.7\%} % all tiers U_fleet
\newcommand{\nAllJ}{7} % all tiers J
\newcommand{\nAllM}{82} % all tiers scenes
\newcommand{\nAllUpop}{15.4\%} % all tiers U_pop
\newcommand{\nPmainOP}{0.88\%} % 1.0000/114 (phase-expected)
\newcommand{\nPmainOPshort}{0.9\%} % same, 1 decimal (abstract)
\newcommand{\nRealMissExpOP}{1.00} % expected real-miss events at OP* (main)
\newcommand{\nRealMissEvOP}{1} % events with a real miss at >= 1 phase
\newcommand{\nRealMissScOP}{1} % scenes with a real miss
\newcommand{\nRecallOP}{100\%} % 1.00/1.00
\newcommand{\nRecallNight}{8\%} % 0.75/9.06
\newcommand{\nChRealMissEv}{13} % (16,16) real-miss events
\newcommand{\nChTwinAlso}{13} % (16,16) twin also misses
\newcommand{\nChRecall}{88\%} % (16,16) recall expected
\newcommand{\nChN}{129} % (16,16) events
\newcommand{\nOldPhat}{11.3\%} % operating_points_v2.final_operating_point.p_hat over 161 CDnet events (all tiers)
\newcommand{\nOldN}{161} % events behind the pooled 11.3%
\newcommand{\nOldNreq}{112} % direct n at 11.3%
\newcommand{\nLosoOP}{0} % LOSO violations D_s > U_pop(-s), OP* (of 48)
\newcommand{\nLosoCh}{0} % LOSO violations, (16,16) (of 48)
\newcommand{\nLosoOPFC}{0} % OP* LOSO false certificates
\newcommand{\nLosoChFC}{0} % (16,16) LOSO false certificates
\newcommand{\nChUpop}{12.5\%} % U_CP(2,48)
\newcommand{\nChJ}{2} % (16,16) J
\newcommand{\nMcFC}{0.0\%} % MC 10,000 false-cert rate OP*
\newcommand{\nScenesPgtEpsOP}{1} % static scenes with p_s > eps at OP* (validity_audit loso.scenes_p_gt_eps)
\newcommand{\nMcFCch}{0.0\%} % MC false-cert rate (16,16)
\newcommand{\nMustFailScenes}{3} % K=48 scenes with p_s > eps
\newcommand{\nMustFailRefuse}{100\%} % refusal rate
\newcommand{\nMustFailApp}{1} % scenes with detector-miss p > eps
\newcommand{\nFleetMeasPB}{0.0272} % max exact PB failure over 12 cases
\newcommand{\nFleetMeasCases}{12} % cases
\newcommand{\nPopTightMax}{0.041} % C4a-pop tight-case max failure (MC 1e5)
\newcommand{\nLeakUpopZero}{6.1\%} % k=0 U_pop OP*
\newcommand{\nLeakUpopSix}{18.5\%} % k=6 U_pop OP*
\newcommand{\nLeakUpopZeroCh}{12.5\%} % k=0 U_pop (16,16)
\newcommand{\nLeakUpopSixCh}{23.0\%} % k=6 (16,16)
\newcommand{\nLeakWithdrawK}{3} % smallest k with all subsets withdrawn (U > eps/2), OP*
\newcommand{\nLeakWithdrawKfirst}{2} % smallest k with some subsets withdrawn, OP*
\newcommand{\nNightFC}{4} % night cameras falsely certified, judged alone (OP*)
\newcommand{\nNightFCch}{3} % same at (16,16)
\newcommand{\nNightPgt}{4} % night cameras with p_s > eps
\newcommand{\nMainFCleak}{0} % main cameras falsely certified, any k
\newcommand{\nSavedOP}{14} % median real events saved per camera, OP*, q known
\newcommand{\nSavedCamOP}{46} % cameras saving (of 48)
\newcommand{\nSavedFinCamOP}{3} % cameras saving with the twin as built
\newcommand{\nSavedFinOP}{0} % median saved, twin as built
\newcommand{\nTwinRunsOP}{28} % exchange_rate.marked.OP*.q_known.n_twin_events_needed_median
\newcommand{\nCpuPerSavedOP}{0.069} % CPU-hours per saved real event, OP* median
\newcommand{\nBreakEven}{14} % 1/c: events/hour (assumes 1 CPU-hour = 1 hour of waiting)
\newcommand{\nSavedCh}{30} % (16,16) saved per camera
\newcommand{\nSavedOPC}{46} % OP-C saved per camera
\newcommand{\nTwinCostScene}{0.20} % CPU-hours per scene, median (C0 build)
\newcommand{\nTwinCostTotal}{14.1} % CPU-hours, 48 main-domain scenes
\newcommand{\nGridPts}{70} % grid points eps=20%
\newcommand{\nGridFinMax}{4} % max cameras saving with the twin as built, any grid point
\newcommand{\nCtwoTwinPass}{73.3\%} % twin pass probability (MC 1e5)
\newcommand{\nCtwoTwinPassShort}{73\%} % twin pass probability, rounded: share of the exact credit realised on average
\newcommand{\nCtwoFCA}{6.7\%} % false cert at p=eps, delta'=0.075
\newcommand{\nCtwoFCB}{8.4\%} % false cert at p=eps, delta'=0.1
\newcommand{\nCtwoFCC}{12.2\%} % false cert at p=eps, delta'=0.15
\newcommand{\nCtwoFCD}{15.6\%} % false cert at p=eps, delta'=0.2
\newcommand{\nCtwoBMCPass}{0\%} % BMC-synth pass probability
\newcommand{\nCtwoBMCN}{30} % BMC-synth events d>=32
\newcommand{\nIWtemporal}{0.64} % iw_ablation_v2 all targets, logistic, temporal
\newcommand{\nIWgeometry}{0.94} % all targets, geometry
\newcommand{\nIWappearance}{0.77} % all targets, appearance only
\newcommand{\nIWtemporalCD}{0.61} % CDnet targets (supplement)
\newcommand{\nIWgeometryCD}{0.96} % CDnet targets
\newcommand{\nIWappearanceCD}{0.75} % CDnet targets
\newcommand{\nClampShare}{70\%} % share of twin frames with a width clamp
\newcommand{\nClampExtraDrop}{3.0} % extra twin hit-rate drop (points) on clamped frames
\newcommand{\nClampPairs}{631} % event x variant pairs
\newcommand{\nCfiveQzero}{1.37\%} % q at x=0
\newcommand{\nCfiveQfifty}{2.36\%} % q at x=50%
\newcommand{\nConfTwin}{0.78} % median detector confidence on twin hits
\newcommand{\nConfReal}{0.87} % median detector confidence on real hits
\newcommand{\nOrigN}{116} % original definition events at OP*
\newcommand{\nOrigJ}{0} % original definition J
\newcommand{\nPmainOPA}{0.88\%} % OP-A main p-hat
\newcommand{\nDexpOPA}{0.00} % OP-A expected discordant events
\newcommand{\nPmainOPC}{0.88\%} % OP-C main p-hat
\newcommand{\nDexpOPC}{0.00} % OP-C expected discordant events
\newcommand{\nTurbDw}{1} % turbulence worst-phase discordant events
\newcommand{\nTurbUfleet}{97.5\%} % turbulence U_fleet
\newcommand{\nTurbJ}{1} % turbulence J
\newcommand{\nTurbM}{2} % turbulence scenes run
\newcommand{\nTurbUpop}{97.5\%} % turbulence U_pop
\newcommand{\nUpopJzero}{6.1\%} % U_CP(0,48), formula (worked example)
\newcommand{\nUpopJone}{9.5\%} % U_CP(1,48), formula (worked example)
\newcommand{\nUpopJtwo}{12.5\%} % U_CP(2,48), formula (worked example)
\newcommand{\nUpopJthree}{15.4\%} % U_CP(3,48), formula (worked example)
\newcommand{\nPiStarA}{6.1\%} % Result 2 threshold pi* on the simulator pass probability, delta'=0.075, eps=delta=5%
\newcommand{\nPiStarB}{3.0\%} % Result 2 threshold pi* on the simulator pass probability, delta'=0.1, eps=delta=5%
\newcommand{\nPiStarC}{1.5\%} % Result 2 threshold pi* on the simulator pass probability, delta'=0.15, eps=delta=5%
\newcommand{\nPiStarD}{1.0\%} % Result 2 threshold pi* on the simulator pass probability, delta'=0.2, eps=delta=5%
\newcommand{\nPAccept}{95.8\%} % share of static cameras accepted by the twin route at OP* (cameras saving / 48)
\newcommand{\nAcceptInflation}{0.3} % U_delta/2(J,48) (1/P(accept) - 1), percentage points
\newcommand{\nUhalfJzero}{7.4\%} % U_CP(0,48) at delta/2 (decision level)
\newcommand{\nUhalfJone}{11.1\%} % U_CP(1,48) at delta/2 (decision level)
\newcommand{\nUhalfJtwo}{14.3\%} % U_CP(2,48) at delta/2 (decision level)
\newcommand{\nUhalfJthree}{17.2\%} % U_CP(3,48) at delta/2 (decision level)
\newcommand{\nPopScenesQuarter}{119} % eps=5%: smallest m with U_CP(0,m) <= eps/2 = 2.5% (J=0, delta=5%)
\newcommand{\nMCScenes}{4} % calibration_scenes.excluded_MC (LASIESTA moving camera)
\newcommand{\nCalibPerCam}{2.4} % annotated calibration events at OP* per static camera (N/m)
\newcommand{\nDirectPooledU}{4.1\%} % direct CP on the calibrated cameras' real misses (events missed at some phase / N)
\newcommand{\nCtwoCFA}{6.7\%} % closed form P(pass)(1-eps)^n1+(1-P(pass))(1-eps)^n0, delta'=0.075
\newcommand{\nCtwoCFB}{8.6\%} % closed form P(pass)(1-eps)^n1+(1-P(pass))(1-eps)^n0, delta'=0.1
\newcommand{\nCtwoCFC}{12.3\%} % closed form P(pass)(1-eps)^n1+(1-P(pass))(1-eps)^n0, delta'=0.15
\newcommand{\nCtwoCFD}{15.5\%} % closed form P(pass)(1-eps)^n1+(1-P(pass))(1-eps)^n0, delta'=0.2
\newcommand{\nCtwoSE}{0.1\%} % MC standard error bound sqrt(0.1*0.9/1e5)=0.09 pts, rounded to 0.1
\newcommand{\nNightHeldD}{23.2\%} % night held-out median per-scene D (pre-registered criterion <= 5%)
\newcommand{\nNightHeldN}{14} % night held-out events at OP*
\newcommand{\nNightBaseD}{47.2\%} % night tuning scenes, base twin, D-hat (phase-averaged)
\newcommand{\nNightProfD}{47.0\%} % night tuning scenes, kept profile, D-hat
\newcommand{\nDefAllA}{161} % primary definition, events of all durations (static)
\newcommand{\nDefAllO}{222} % original definition, events of all durations (static)
\newcommand{\nOrigChJ}{3} % original definition (16,16) J
\newcommand{\nOrigChUpop}{15.4\%} % original (16,16) U_pop
\newcommand{\nOrigOPUfleet}{2.55\%} % original OP* U_fleet (worst phase)
\newcommand{\nOrigOPDw}{0} % original OP* disc
\newcommand{\nTpOPS}{0.88\%} % OP* static-domain p-hat
\newcommand{\nTndOPS}{14} % OP* n_direct (80% power at p-hat)
\newcommand{\nTsvOPS}{14} % OP* saved per camera (median)
\newcommand{\nTcamOPS}{46} % OP* cameras saving (of 48)
\newcommand{\nTcamfinOPS}{3} % OP* cameras saving, twin as built
\newcommand{\nTrunsOPS}{28} % OP* twin runs needed per camera
\newcommand{\nTcpuOPS}{0.069} % OP* CPU-h per saved event
\newcommand{\nTpOPA}{0.88\%} % OP-A static-domain p-hat
\newcommand{\nTndOPA}{14} % OP-A n_direct (80% power at p-hat)
\newcommand{\nTsvOPA}{14} % OP-A saved per camera (median)
\newcommand{\nTcamOPA}{47} % OP-A cameras saving (of 48)
\newcommand{\nTcamfinOPA}{3} % OP-A cameras saving, twin as built
\newcommand{\nTrunsOPA}{28} % OP-A twin runs needed per camera
\newcommand{\nTcpuOPA}{0.064} % OP-A CPU-h per saved event
\newcommand{\nTpOPC}{0.88\%} % OP-C static-domain p-hat
\newcommand{\nTndOPC}{46} % OP-C n_direct (80% power at p-hat)
\newcommand{\nTsvOPC}{46} % OP-C saved per camera (median)
\newcommand{\nTcamOPC}{47} % OP-C cameras saving (of 48)
\newcommand{\nTcamfinOPC}{0} % OP-C cameras saving, twin as built
\newcommand{\nTrunsOPC}{149} % OP-C twin runs needed per camera
\newcommand{\nTcpuOPC}{0.104} % OP-C CPU-h per saved event
\newcommand{\nTpCH}{4.26\%} % challenge_16_16 static-domain p-hat
\newcommand{\nTndCH}{30} % challenge_16_16 n_direct (80% power at p-hat)
\newcommand{\nTsvCH}{30} % challenge_16_16 saved per camera (median)
\newcommand{\nTcamCH}{45} % challenge_16_16 cameras saving (of 48)
\newcommand{\nTcamfinCH}{0} % challenge_16_16 cameras saving, twin as built
\newcommand{\nTrunsCH}{66} % challenge_16_16 twin runs needed per camera
\newcommand{\nTcpuCH}{0.078} % challenge_16_16 CPU-h per saved event
\newcommand{\nWPSpMain}{0.9\%} % worst-phase real miss rate, main, OP* (events missed at >= 1 phase / N=114)
\newcommand{\nWPSqMain}{3.5\%} % worst-phase twin miss rate (majority), main, OP*
\newcommand{\nPASpMain}{0.9\%} % phase-averaged real miss rate, main, OP*
\newcommand{\nPASqMain}{1.4\%} % phase-averaged twin miss rate (majority), main, OP*
\newcommand{\nWPSpJit}{4.8\%} % worst-phase real miss rate, jitter, OP* (events missed at >= 1 phase / N=62)
\newcommand{\nWPSqJit}{0.0\%} % worst-phase twin miss rate (majority), jitter, OP*
\newcommand{\nPASpJit}{0.5\%} % phase-averaged real miss rate, jitter, OP*
\newcommand{\nPASqJit}{0.0\%} % phase-averaged twin miss rate (majority), jitter, OP*
\newcommand{\nWPSpNight}{52.0\%} % worst-phase real miss rate, night, OP* (events missed at >= 1 phase / N=25)
\newcommand{\nWPSqNight}{16.0\%} % worst-phase twin miss rate (majority), night, OP*
\newcommand{\nPASpNight}{36.2\%} % phase-averaged real miss rate, night, OP*
\newcommand{\nPASqNight}{4.8\%} % phase-averaged twin miss rate (majority), night, OP*
\newcommand{\nWPSpTurb}{50.0\%} % worst-phase real miss rate, turbulence, OP* (events missed at >= 1 phase / N=2)
\newcommand{\nWPSqTurb}{0.0\%} % worst-phase twin miss rate (majority), turbulence, OP*
\newcommand{\nPASpTurb}{6.2\%} % phase-averaged real miss rate, turbulence, OP*
\newcommand{\nPASqTurb}{0.0\%} % phase-averaged twin miss rate (majority), turbulence, OP*
\newcommand{\nWPSpAll}{8.9\%} % worst-phase real miss rate, all_tiers, OP* (events missed at >= 1 phase / N=203)
\newcommand{\nWPSqAll}{3.9\%} % worst-phase twin miss rate (majority), all_tiers, OP*
\newcommand{\nPASpAll}{5.2\%} % phase-averaged real miss rate, all_tiers, OP*
\newcommand{\nPASqAll}{1.4\%} % phase-averaged twin miss rate (majority), all_tiers, OP*
\newcommand{\nWXSp}{0.9\%} % pooled worst-phase real miss rate, main, OP*
\newcommand{\nWXSnd}{14} % direct plan at the worst-phase p, OP*
\newcommand{\nWXSacc}{44} % cameras accepted with worst-phase q (of 48), OP*
\newcommand{\nWXSsaved}{14} % median real events saved per camera, worst-phase, OP*
\newcommand{\nWXSfa}{0} % accepted cameras with worst-phase p_s > eps, OP*
\newcommand{\nWXSruns}{28} % twin runs needed per camera (median), worst-phase q, OP*
\newcommand{\nWPSDwMain}{0} % worst-phase discordant events, main, OP* (of N=114)
\newcommand{\nWPSNMain}{114} % main events, OP*
\newcommand{\nWPSUfMain}{2.6\%} % U_CP(0,114), main, OP*
\newcommand{\nWPSJMain}{0} % main scenes with a worst-phase discordant event, OP* (of 48)
\newcommand{\nWPSUpMain}{6.1\%} % U_CP(0,48), main, OP*
\newcommand{\nWPSNJit}{62} % jitter events, OP*
\newcommand{\nWPSDwJit}{3} % jitter worst-phase discordant events, OP*
\newcommand{\nOODSJitcams}{26} % jitter cameras judged with a main-only calibration, OP*
\newcommand{\nOODSJitacc}{26} % jitter cameras accepted (worst-phase q_w + U_delta/2(J_main,m_main) <= eps), OP*
\newcommand{\nOODSJitfa}{1} % jitter cameras falsely accepted (accepted and p_w > eps), OP*
\newcommand{\nOODSJitpgt}{1} % jitter cameras with p_w > eps, OP*
\newcommand{\nOODSNightcams}{6} % night cameras judged with a main-only calibration, OP*
\newcommand{\nOODSNightacc}{5} % night cameras accepted (worst-phase q_w + U_delta/2(J_main,m_main) <= eps), OP*
\newcommand{\nOODSNightfa}{4} % night cameras falsely accepted (accepted and p_w > eps), OP*
\newcommand{\nOODSNightpgt}{5} % night cameras with p_w > eps, OP*
\newcommand{\nOODSJitFApw}{37.5\%} % p_w of the falsely accepted jitter camera traffic, OP*
\newcommand{\nWPCpMain}{10.1\%} % worst-phase real miss rate, main, challenge_16_16 (events missed at >= 1 phase / N=129)
\newcommand{\nWPCqMain}{13.2\%} % worst-phase twin miss rate (majority), main, challenge_16_16
\newcommand{\nPACpMain}{4.3\%} % phase-averaged real miss rate, main, challenge_16_16
\newcommand{\nPACqMain}{6.5\%} % phase-averaged twin miss rate (majority), main, challenge_16_16
\newcommand{\nWPCpJit}{10.4\%} % worst-phase real miss rate, jitter, challenge_16_16 (events missed at >= 1 phase / N=67)
\newcommand{\nWPCqJit}{0.0\%} % worst-phase twin miss rate (majority), jitter, challenge_16_16
\newcommand{\nPACpJit}{1.9\%} % phase-averaged real miss rate, jitter, challenge_16_16
\newcommand{\nPACqJit}{0.0\%} % phase-averaged twin miss rate (majority), jitter, challenge_16_16
\newcommand{\nWPCpNight}{63.9\%} % worst-phase real miss rate, night, challenge_16_16 (events missed at >= 1 phase / N=36)
\newcommand{\nWPCqNight}{30.6\%} % worst-phase twin miss rate (majority), night, challenge_16_16
\newcommand{\nPACpNight}{49.7\%} % phase-averaged real miss rate, night, challenge_16_16
\newcommand{\nPACqNight}{8.0\%} % phase-averaged twin miss rate (majority), night, challenge_16_16
\newcommand{\nWPCpTurb}{50.0\%} % worst-phase real miss rate, turbulence, challenge_16_16 (events missed at >= 1 phase / N=2)
\newcommand{\nWPCqTurb}{0.0\%} % worst-phase twin miss rate (majority), turbulence, challenge_16_16
\newcommand{\nPACpTurb}{6.2\%} % phase-averaged real miss rate, turbulence, challenge_16_16
\newcommand{\nPACqTurb}{0.0\%} % phase-averaged twin miss rate (majority), turbulence, challenge_16_16
\newcommand{\nWPCpAll}{18.8\%} % worst-phase real miss rate, all_tiers, challenge_16_16 (events missed at >= 1 phase / N=234)
\newcommand{\nWPCqAll}{12.0\%} % worst-phase twin miss rate (majority), all_tiers, challenge_16_16
\newcommand{\nPACpAll}{10.6\%} % phase-averaged real miss rate, all_tiers, challenge_16_16
\newcommand{\nPACqAll}{4.8\%} % phase-averaged twin miss rate (majority), all_tiers, challenge_16_16
\newcommand{\nWXCp}{10.1\%} % pooled worst-phase real miss rate, main, challenge_16_16
\newcommand{\nWXCnd}{88} % direct plan at the worst-phase p, challenge_16_16
\newcommand{\nWXCacc}{43} % cameras accepted with worst-phase q (of 48), challenge_16_16
\newcommand{\nWXCsaved}{88} % median real events saved per camera, worst-phase, challenge_16_16
\newcommand{\nWXCfa}{0} % accepted cameras with worst-phase p_s > eps, challenge_16_16
\newcommand{\nWXCruns}{66} % twin runs needed per camera (median), worst-phase q, challenge_16_16
\newcommand{\nWPCDwMain}{5} % worst-phase discordant events, main, challenge_16_16 (of N=129)
\newcommand{\nWPCNMain}{129} % main events, challenge_16_16
\newcommand{\nWPCUfMain}{8.0\%} % U_CP(5,129), main, challenge_16_16
\newcommand{\nWPCJMain}{2} % main scenes with a worst-phase discordant event, challenge_16_16 (of 48)
\newcommand{\nWPCUpMain}{12.5\%} % U_CP(2,48), main, challenge_16_16
\newcommand{\nWPCNJit}{67} % jitter events, challenge_16_16
\newcommand{\nWPCDwJit}{7} % jitter worst-phase discordant events, challenge_16_16
\newcommand{\nOODCJitcams}{26} % jitter cameras judged with a main-only calibration, challenge_16_16
\newcommand{\nOODCJitacc}{26} % jitter cameras accepted (worst-phase q_w + U_delta/2(J_main,m_main) <= eps), challenge_16_16
\newcommand{\nOODCJitfa}{1} % jitter cameras falsely accepted (accepted and p_w > eps), challenge_16_16
\newcommand{\nOODCJitpgt}{1} % jitter cameras with p_w > eps, challenge_16_16
\newcommand{\nOODCNightcams}{6} % night cameras judged with a main-only calibration, challenge_16_16
\newcommand{\nOODCNightacc}{3} % night cameras accepted (worst-phase q_w + U_delta/2(J_main,m_main) <= eps), challenge_16_16
\newcommand{\nOODCNightfa}{3} % night cameras falsely accepted (accepted and p_w > eps), challenge_16_16
\newcommand{\nOODCNightpgt}{6} % night cameras with p_w > eps, challenge_16_16
\newcommand{\nOODCJitFApw}{53.8\%} % p_w of the falsely accepted jitter camera traffic, challenge_16_16
\newcommand{\nSiteSm}{17} % sites (CDnet categories + LASIESTA groups), OP*
\newcommand{\nSiteSJ}{0} % sites with a discordant event, OP*
\newcommand{\nSiteSU}{16.2\%} % site-level population bound U_CP(J_site, m_site), OP*
\newcommand{\nSiteSloso}{0} % leave-one-site-out violations, OP*
\newcommand{\nSiteCm}{17} % sites (CDnet categories + LASIESTA groups), challenge_16_16
\newcommand{\nSiteCJ}{1} % sites with a discordant event, challenge_16_16
\newcommand{\nSiteCU}{25.0\%} % site-level population bound U_CP(J_site, m_site), challenge_16_16
\newcommand{\nSiteCloso}{1} % leave-one-site-out violations, challenge_16_16
\newcommand{\nRepoTag}{ress-v1.2} % artifact release tag
\newcommand{\nRepoHashShort}{890465bd9267} % commit A (tagged) short hash

% generated by tools/make_proofs_numbers.py from results/*.json -- do not edit
\newcommand{\numConeClosed}{0.0485}
\newcommand{\numConeMC}{0.0491}
\newcommand{\numCtwoMaxDev}{0.29}
\newcommand{\numCtwoNgrid}{40}
\newcommand{\numCtwoAllMax}{yes}
\newcommand{\numCtwoTwinPass}{98.2\%}
\newcommand{\numCtwoFCA}{7.2\%}
\newcommand{\numCtwoCredA}{7.9}
\newcommand{\numCtwoFCB}{10.1\%}
\newcommand{\numCtwoCredB}{13.7}
\newcommand{\numCtwoFCC}{14.8\%}
\newcommand{\numCtwoCredC}{21.6}
\newcommand{\numCtwoFCD}{19.3\%}
\newcommand{\numCtwoCredD}{26.5}
\newcommand{\numCtwoBMCPass}{0.0\%}
\newcommand{\numCthreeA}{72}
\newcommand{\numCthreeB}{91}
\newcommand{\numCthreeC}{122}
\newcommand{\numLemmaMin}{0.2500}
\newcommand{\numFleetPBmax}{0.0471}
\newcommand{\numFleetMCmax}{0.0473}
\newcommand{\numFleetNpat}{6}
\newcommand{\numFleetMeasPBmax}{0.0126}
\newcommand{\numFleetMeasMCmax}{0.0129}
\newcommand{\numFleetMeasN}{15}
\newcommand{\numPopFailMax}{0.0401}
\newcommand{\numCfourbFailMax}{0.0286}
\newcommand{\numPopScenesTen}{29}
\newcommand{\numPopScenesFive}{59}
\newcommand{\numLosoViolNum}{2}
\newcommand{\numLosoViolDen}{78}
\newcommand{\numCfiveAll}{yes}
\newcommand{\numCfiveQzero}{0.87\%}
\newcommand{\numCfiveQfifty}{1.67\%}


\newtheorem{theorem}{Theorem}
\newtheorem{proposition}[theorem]{Proposition}
\newtheorem{lemma}[theorem]{Lemma}
\newtheorem{corollary}[theorem]{Corollary}
\theoremstyle{definition}
\newtheorem{definition}{Definition}
\theoremstyle{remark}
\newtheorem{remark}{Remark}
\newcommand{\Bin}{\mathrm{Bin}}
\newcommand{\Bern}{\mathrm{Bern}}
\newcommand{\PP}{\mathbb{P}}
\newcommand{\EE}{\mathbb{E}}
\newcommand{\UCP}{U}
\renewcommand{\thesection}{S\arabic{section}}
\renewcommand{\thetable}{S\arabic{table}}
\floatplacement{table}{H}

\title{Supplementary material for\\ ``No Free Credit: How Much Simulated Evidence Can Replace Real Events in
Zero-Failure Reliability Demonstration of Perception Systems''}
\author{Lam Quoc Dat\\ \small School of Business and Technology (FSB), FPT University, Ho Chi Minh City, Vietnam\\ \small lamquocdat@gmail.com, ORCID 0009-0004-5432-9343}
\date{}

\begin{document}
\maketitle

\noindent Contents: S1 proofs of C1--C5a with numerical checks; S2 units and event counts; S3 the full operating-point
grid; S4 audit details and must-fail scenarios; S5 sensitivity analyses; S6 twin construction and asset licences;
S7 all supplementary tables.
All numbers are generated from the study's result files; Monte-Carlo checks use seed 42.

% ------------------------------------------------------------------------------------------------------
\section{Proofs of C1--C5a}\label{sec:proofs}
Numbering in this section is local. Correspondence with the main text: Theorem~S1 proves Result~1, Theorem~S2
Result~2, Lemma~S3 and Proposition~S9 Result~3, Theorem~S6, Proposition~S7 and Proposition~S8 Result~4(a), (b)
and (c). The main text states the same results in its Appendix~A.
\renewcommand{\thetheorem}{S\arabic{theorem}}
{\let\section\subsection
% body of the proofs; \input by proofs_v1.tex and manuscript/supplement.tex
\section{Setting and notation}\label{sec:setting}

\paragraph{Events and misses.} A camera (scene) $s$ produces events $i=1,2,\dots$ (a maximal run of frames containing
the target object). The system samples frames periodically with period $K$ and an unknown phase
$\varphi\in\{0,\dots,K-1\}$; a detector is run on the sampled frames. An event is \emph{missed} ($M=1$) when no
sampled frame of the event is a hit (the detector finds the object). With $H\subseteq\mathbb{Z}$ the hit frames of
the event and $\varphi$ uniform, $\PP(M=1\mid H)=1-|H \bmod K|/K$, where $H\bmod K$ is the set of residues of $H$.
The miss rate of camera $s$ is $p_s=\PP(M=1)$ over its events and phases. Target: certify $p_s\le\varepsilon$ with
confidence $1-\delta$.

\paragraph{Clopper--Pearson bound.} For $x\in\{0,\dots,n\}$, $F_{n,p}(x)=\PP(\Bin(n,p)\le x)$ and
$\UCP(x,n)=\UCP_\delta(x,n)$ is the one-sided $(1-\delta)$ upper Clopper--Pearson bound
\cite{clopper1934}: the solution $u$ of $F_{n,u}(x)=\delta$ for $x<n$, and $1$ for $x=n$. Since $F_{n,p}(x)$ is
strictly decreasing in $p$, $\UCP(x,n)<p \iff F_{n,p}(x)<\delta$, and $\UCP$ is non-decreasing in $x$. Validity:
$\PP(\UCP(\Bin(n,p),n)<p)\le\delta$ for every $p$.

\paragraph{Success run.} $n_0(\varepsilon,\delta)=\lceil \ln\delta/\ln(1-\varepsilon)\rceil$ is the smallest $n$ with
$(1-\varepsilon)^n\le\delta$ (e.g.\ $n_0(5\%,5\%)=59$, $n_0(20\%,5\%)=14$).

\paragraph{Simulator and twin.} Simulator data $Y$ have a law $Q$ on an arbitrary space. A \emph{twin} is a paired
simulator: every real event $i$ has a synthetic counterpart with miss indicator $M'_i$ under the same sampling phase.
$q=\PP(M'=1)$, and the \emph{discordance} is $D=\PP(M=1,M'=0)$ (real miss, twin catch).

\section{C1: no free credit}

A certification procedure is a (possibly randomised) test $\phi(X,Y)\in\{0,1\}$ with real data
$X=(M_1,\dots,M_n)\sim\Bern(p)^{\otimes n}$ independent of simulator data $Y\sim Q$. It is \emph{valid at level
$\delta$ on a class $\mathcal{P}$} of pairs $(p,Q)$ if $\sup_{(p,Q)\in\mathcal{P},\,p>\varepsilon}\PP_{p,Q}(\phi=1)\le\delta$.

\begin{theorem}[C1]\label{thm:c1}
Assume $\mathcal{P}$ places no restriction linking $Q$ and $p$: for every $Q$ there are $(0,Q)\in\mathcal{P}$ and
$(p,Q)\in\mathcal{P}$ for all $p\in(\varepsilon,\varepsilon+\eta)$ for some $\eta>0$. If $\phi$ is valid at level
$\delta$ on $\mathcal{P}$, then for every $Q$ and every simulator sample size,
\[(1-\varepsilon)^n\,\PP_{0,Q}(\phi=1)\le\delta .\]
In particular, certifying a perfect system ($p=0$) with probability one requires $n\ge n_0(\varepsilon,\delta)$ real
events, whatever the amount of simulator data; power $1-\beta$ at $p=0$ requires
$n\ge \ln(\delta/(1-\beta))/\ln(1-\varepsilon)$.
\end{theorem}

\begin{proof}
Fix $Q$ and write $g(y)=\EE[\phi(0^n,y)]$ (expectation over the internal randomisation). Under $p=0$, $X=0^n$ almost
surely, so $\PP_{0,Q}(\phi=1)=\EE_Q[g(Y)]$. Under $p>\varepsilon$, the event $\{X=0^n\}$ has probability $(1-p)^n$ and
is independent of $Y$, hence $\PP_{p,Q}(\phi=1)\ge(1-p)^n\,\EE_Q[g(Y)]=(1-p)^n\PP_{0,Q}(\phi=1)$. Validity gives
$(1-p)^n\PP_{0,Q}(\phi=1)\le\delta$ for all $p\in(\varepsilon,\varepsilon+\eta)$; letting $p\downarrow\varepsilon$
proves the inequality. With $\PP_{0,Q}(\phi=1)=1$ it reads $(1-\varepsilon)^n\le\delta$, i.e.\ $n\ge n_0$; with
$\PP_{0,Q}(\phi=1)\ge 1-\beta$ it gives the power bound.
\end{proof}

\begin{remark}
The argument only uses that the simulator data have the same law under the two hypotheses and that the
all-success real sample has positive probability under both. It is the success-run analogue of the observation that
external data cannot increase power under strict type-I error control over all external-data laws
\cite{koppschneider2020}.
\emph{Check.} The zero-failure test with $n=59$ at $p=\varepsilon=5\%$: $(0.95)^{59}=\numConeClosed$ (MC
$\numConeMC$, $10^5$ runs).
\end{remark}

\section{C2: the price of borrowing}

Suppose one accepts a two-tier guarantee: level $\delta$ when the simulator is ``right'' and only level
$\delta'>\delta$ in general.

\begin{theorem}[C2, exact credit]\label{thm:c2}
(i) \emph{Converse.} Under the assumption of Theorem~\ref{thm:c1}, any procedure valid at level $\delta'$ on
$\mathcal{P}$ that certifies a perfect system with probability one uses $n\ge n_0(\varepsilon,\delta')$ real events.
Hence the credit, the number of real events saved with respect to $n_0(\varepsilon,\delta)$, is at most
\[c=n_0(\varepsilon,\delta)-n_0(\varepsilon,\delta').\]
(ii) \emph{Attainment.} Let $T(Y)\in\{0,1\}$ be any simulator test (``the simulator passes'') and let $B$ be the
procedure: if $T=1$ run the zero-failure test with $n_1=n_0(\varepsilon,\delta')$ real events, otherwise with
$n_0=n_0(\varepsilon,\delta)$. Then $B$ is valid at level $\delta'$ on $\mathcal{P}$, saves exactly $c$ events
whenever $T=1$, and is valid at level $\delta$ on the subclass
\[\mathcal{P}_0=\Bigl\{(p,Q):\ p>\varepsilon\Rightarrow \PP_Q(T=1)\le\pi^\ast\Bigr\},\qquad
\pi^\ast=\frac{\delta-(1-\varepsilon)^{n_0}}{(1-\varepsilon)^{n_1}-(1-\varepsilon)^{n_0}},\]
i.e.\ where a simulator of a system that is truly worse than $\varepsilon$ rarely passes.
(iii) \emph{Continuous lower bound.} $c\ \ge\ c^\ast=\bigl\lfloor \ln(\delta'/\delta)/(-\ln(1-\varepsilon))\bigr\rfloor$.
\end{theorem}

\begin{proof}
(i) Theorem~\ref{thm:c1} at level $\delta'$ gives $(1-\varepsilon)^n\le\delta'$, whose smallest integer solution is
$n_0(\varepsilon,\delta')$. (ii) For $p>\varepsilon$, with $\pi=\PP_Q(T=1)$ and $T$ independent of the real data,
$\PP(B=1)=\pi(1-p)^{n_1}+(1-\pi)(1-p)^{n_0}\le(1-\varepsilon)^{n_1}\le\delta'$, and $\le\delta$ as soon as
$\pi\le\pi^\ast$ (the expression is affine in $\pi$ and decreasing in $p$). (iii) Write $a=\ln\delta/\ln(1-\varepsilon)$,
$a'=\ln\delta'/\ln(1-\varepsilon)$ and $r=a-a'=\ln(\delta'/\delta)/(-\ln(1-\varepsilon))$. Since $n_0(\delta)\ge a$ and
$n_0(\delta')<a'+1$, $c>r-1$; $c$ being an integer, $c\ge\lfloor r\rfloor=c^\ast$.
\end{proof}

\begin{remark}[numbers and checks]
At $\varepsilon=\delta=5\%$: $\delta'=7.5,10,15,20\%$ give $c=8,14,22,27$ ($-14\%,-24\%,-37\%,-46\%$ of 59).
The exact form is maximal on all \numCtwoNgrid{} grid points ($(1-\varepsilon)^{n_0(\delta')}\le\delta'<(1-\varepsilon)^{n_0(\delta')-1}$:
\numCtwoAllMax); MC ($10^5$ runs per point) matches the closed form within \numCtwoMaxDev{} percentage points.
\emph{Empirical C2} (procedure $B$, $T$ = the one-sided CP bound of the simulator's miss rate $\le\varepsilon$,
$10^5$ runs, $\varepsilon=5\%$, $d_{\min}=32$, $K=1$): with the paired twin of 78 static scenes (180 events) the
simulator passes in \numCtwoTwinPass{} of runs; realised credit $\numCtwoCredA/\numCtwoCredB/\numCtwoCredC/\numCtwoCredD$
events and false certification at the worst case $p=\varepsilon$ equal to
$\numCtwoFCA/\numCtwoFCB/\numCtwoFCC/\numCtwoFCD$ for $\delta'=7.5/10/15/20\%$ --- the credit is bought exactly with
the extra error $\delta'-\delta$. The off-the-shelf synthetic benchmark (30 events of length $\ge 32$) never passes
(\numCtwoBMCPass): too few synthetic events to bound its own miss rate below $5\%$.
\end{remark}

\section{C3: the twin lemma}

\begin{lemma}[C3]\label{lem:c3}
For any pairing $(M,M')$ of a real event and its twin, $p\le q+D$ with $D=\PP(M=1,M'=0)$. No assumption links the
twin to reality.
\end{lemma}

\begin{proof}
$p=\PP(M=1,M'=1)+\PP(M=1,M'=0)\le\PP(M'=1)+D=q+D$.
\end{proof}

\begin{corollary}[a per-camera twin does not save events]
If $q\le q_U$ holds with confidence $1-\delta/2$ and $D\le\varepsilon-q_U$ is certified for the same camera by a
zero-discordance success run at level $\delta/2$, the number of paired real events is
$\lceil\ln(\delta/2)/\ln(1-(\varepsilon-q_U))\rceil$ $=\numCthreeA,\numCthreeB,\numCthreeC$ for $q_U=0,1,2\%$ at
$\varepsilon=\delta=5\%$, all larger than $n_0=59$. The twin can only help when $D$ is calibrated once, on many
scenes, and transferred (C4).
\end{corollary}

\begin{remark}[event unit]
With a uniform phase, $x_i=\PP_\varphi(M_i=1,M'_i=0)=|R'_i\setminus R_i|/K$, where $R_i$, $R'_i$ are the residue
sets of the real and twin hit frames of event $i$ (twin catch decided by the majority of three sprite variants).
The \emph{worst-phase} indicator $W_i=\mathbf{1}\{R'_i\setminus R_i\neq\emptyset\}$ satisfies
$\PP(M_i=1,M'_i=0\mid\text{event})=x_i\le W_i$. All certificates below are stated for Bernoulli indicators
$Z_i$; they hold verbatim with $Z_i=W_i$ and then bound the (larger) worst-phase discordance rate. This integer
form is the one used for the reported certificates.
\end{remark}

\section{fleet bound: event-weighted certificate for the calibrated fleet}

\begin{lemma}[mean--median]\label{lem:mm}
For $N\ge 2$ and $0<p<1-1/N$, $F_{N,p}(\lfloor Np\rfloor)>1/4$. Consequently, if $\delta\le 1/4$ and
$F_{N,p}(x)<\delta$, then $x\le Np-1$.
\end{lemma}

\begin{proof}
Greenberg and Mohri \cite{greenberg2014} prove $\PP(B\ge\EE B)>1/4$ for $B\sim\Bin(N,r)$ with $r>1/N$. Apply it to
$B=N-\Bin(N,p)\sim\Bin(N,1-p)$ with $1-p>1/N$: $\PP(\Bin(N,p)\le Np)>1/4$, and $\{\Bin(N,p)\le Np\}=\{\Bin(N,p)\le\lfloor Np\rfloor\}$.
For the consequence: $F_{N,p}$ is non-decreasing, so $F_{N,p}(x)<\delta\le1/4<F_{N,p}(\lfloor Np\rfloor)$ forces
$x<\lfloor Np\rfloor$, i.e.\ $x\le\lfloor Np\rfloor-1\le Np-1$.
\emph{Check:} $\inf F_{N,p}(\lfloor Np\rfloor)$ over $N=2..2000$ and a fine grid of $p$ is \numLemmaMin, approached at
$N=2$, $p\to1/2$ (the constant $1/4$ is tight).
\end{proof}

\begin{theorem}[fleet bound]\label{thm:fleet}
Let $Z_1,\dots,Z_N$ be independent, $Z_i\sim\Bern(d_i)$ with arbitrary $d_i\in[0,1]$ (heterogeneous across scenes and
events), $X=\sum_iZ_i$ and $\bar d=\frac1N\sum_id_i$. If $\delta\le1/4$ and $\bar d<1-1/N$, then
\[\PP\bigl(\UCP(X,N)<\bar d\bigr)\le\delta .\]
With $d_i=D_{s(i)}$ (the discordance of the scene of event $i$), $\bar d=\bar D_w=\sum_sn_sD_s/N$, the
event-weighted discordance of the calibrated cameras; combined with Lemma~\ref{lem:c3},
$\bar p_w\le\bar q_w+\UCP(X,N)$ with probability at least $1-\delta$.
\end{theorem}

\begin{proof}
If $\bar d=0$ there is nothing to prove. Otherwise let $x^\ast=\max\{x:F_{N,\bar d}(x)<\delta\}$ (if the set is empty
the failure probability is $0$). By monotonicity of $\UCP$, $\{\UCP(X,N)<\bar d\}=\{X\le x^\ast\}$. By
Lemma~\ref{lem:mm}, $x^\ast\le N\bar d-1$. Hoeffding's Theorem~4 \cite{hoeffding1956}: for a sum $S$ of independent
Bernoulli variables with mean $N\bar d$ and $0\le b\le N\bar d-1$, $\PP(S\le b)\le F_{N,\bar d}(b)$. With $b=x^\ast$,
$\PP(X\le x^\ast)\le F_{N,\bar d}(x^\ast)<\delta$.
\end{proof}

\begin{remark}[what it does not say; checks]
The guarantee is conditional on the calibrated scenes and weighted by events: it does not bound any single camera
and does not transfer to a new camera (if scenes are random, $X$ is over-dispersed; e.g.\ $D_s\in\{0,1\}$ gives
$\PP(X=0)=(1-\mu)^m\gg(1-\mu)^N$).
\emph{Checks.} Six synthetic heterogeneity patterns on a fixed configuration of 44 scenes with $N=304$ events (a
calibration set from an earlier stage of the study, used only as a test bed for the theorem): exact
Poisson--binomial failure $\le\numFleetPBmax$, MC ($10^5$) $\le\numFleetMCmax$. Measured discordance probabilities of
the study (event level, main domain, five operating points, as measured and rescaled to mean $2\%$ and $5\%$;
\numFleetMeasN{} cases): exact failure $\le\numFleetMeasPBmax$, MC $\le\numFleetMeasMCmax$, all $\le\delta=0.05$.
\end{remark}

\section{population bound: transfer to a new camera}

\begin{proposition}[population bound]\label{prop:pop}
Let scenes $S_1,\dots,S_m$ be i.i.d.\ from a scene population, each with $n_s\ge1$ paired events whose
discordance indicators are independent $\Bern(D_{S_s})$ given the scene, $\mu=\EE[D_S]$, and
$J=\#\{s:X_s\ge1\}$ with $X_s$ the number of discordant events in scene $s$. Then
\[\PP\bigl(\UCP(J,m)<\mu\bigr)\le\delta .\]
For a new camera drawn from the same population, $\EE[p_{\mathrm{new}}]\le\EE[q_{\mathrm{new}}]+\UCP(J,m)$ with
probability $\ge1-\delta$ over the calibration, and $\PP(D_{\mathrm{new}}>t)\le\UCP(J,m)/t$ (Markov).
\end{proposition}

\begin{proof}
Let $B_s=\mathbf 1\{X_s\ge1\}$. Given the scene, $\PP(B_s=1\mid S_s)=1-(1-D_{S_s})^{n_s}\ge D_{S_s}$ because
$(1-D)^n\le 1-D$ for $n\ge1$. Hence $\beta_s:=\PP(B_s=1)\ge\mu$, and the $B_s$ are independent across scenes. By a
coupling ($B_s=\mathbf 1\{V_s\le\beta_s\}\ge\mathbf 1\{V_s\le\mu\}$ with i.i.d.\ uniforms $V_s$), $J$ stochastically
dominates $\Bin(m,\mu)$. Since $\UCP(\cdot,m)$ is non-decreasing,
$\PP(\UCP(J,m)<\mu)\le\PP(\UCP(\Bin(m,\mu),m)<\mu)\le\delta$ by Clopper--Pearson validity. The corollary follows from
Lemma~\ref{lem:c3} applied to the new camera and averaged over the population.
\end{proof}

\begin{remark}[price and checks]
The resolution is set by the number of scenes: with $J=0$, $\UCP\le10\%$ needs $m\ge\numPopScenesTen$ and
$\UCP\le5\%$ needs $m\ge\numPopScenesFive$. \emph{Checks.} The tight case $D_S\in\{0,1\}$, $n_s=1$, $m=78$ (then
$J\sim\Bin(m,\mu)$ exactly), $\mu\in\{1,2,5,10,20\}\%$, $10^5$ runs: maximum failure \numPopFailMax. Leave-one-scene-out
on the 78 static scenes at the post-hoc challenge point ($d_{\min}=16$, $K=16$): $D_s>\UCP(J_{-s},m-1)$ for
\numLosoViolNum{} of \numLosoViolDen{} scenes (the proposition bounds the population mean, not each scene).
\end{remark}

\section{per-camera form: a per-camera (quantile) certificate and its data requirement}

\begin{proposition}[per-camera form]\label{prop:c4b}
Under the assumptions of Proposition~\ref{prop:pop} with $n_s\ge n_{\min}$ for all $s$, let $t\in(0,1)$,
$\pi_t=\PP(D_S>t)$ and $h_t=1-(1-t)^{n_{\min}}$. Then
\[\PP\bigl(\pi_t>\UCP(J,m)/h_t\bigr)\le\delta .\]
Hence, with $\hat\gamma=\UCP(J,m)/h_t$, a new camera satisfies $p_{\mathrm{new}}\le q_{\mathrm{new}}+t$ except on a
set of cameras of population mass at most $\hat\gamma$, with confidence $1-\delta$.
\end{proposition}

\begin{proof}
$\PP(B_s=1)\ge\PP(D_S>t)\bigl(1-(1-t)^{n_s}\bigr)\ge\pi_th_t$. The coupling argument of Proposition~\ref{prop:pop}
gives $J\succeq_{\mathrm{st}}\Bin(m,\pi_th_t)$, so $\PP(\UCP(J,m)<\pi_th_t)\le\delta$, which is the claim.
\end{proof}

\begin{remark}[data requirement; checks]
With $J=0$, $\hat\gamma\le\gamma$ requires $m\ge m_{\min}$ with $1-\delta^{1/m_{\min}}\le\gamma$ ($m_{\min}=59$ for
$\gamma=5\%$, $29$ for $\gamma=10\%$) and $n_{\min}\ge\ln\bigl(1-\UCP(0,m)/\gamma\bigr)/\ln(1-t)$ events in every scene;
the total $m\,n_{\min}$ decreases to the floor $\approx\ln(1/\delta)/(\gamma t)$ as $m$ grows. \emph{Check:} two-point
populations $D_S\in\{0,t'\}$, $t'>t$, $m=78$, $10^5$ runs: maximum failure \numCfourbFailMax.
\end{remark}

\section{C5a: a pessimistic-length twin is conservative under periodic sampling}

\begin{proposition}[C5a]\label{prop:c5a}
Let an event occupy frames $F=\{o,\dots,o+d-1\}$ and let the detector outcome on each frame be fixed (a deterministic
function of the frame). For the truncated event $F_L=\{o,\dots,o+L-1\}$, $L\le d$, with hit set $H_L=H\cap F_L$, the
miss probability under a uniform phase, $m(L)=1-|H_L\bmod K|/K$, is non-increasing in $L$. Consequently a twin built
with a pessimistic (shorter) length $L$ has $q_L\ge q$ and discordance $D_L\le D$, so the C3 bound $q_L+D_L$ remains
valid and the miss term is conservative.
\end{proposition}

\begin{proof}
$L\le L'$ implies $H_L\subseteq H_{L'}$, hence $H_L\bmod K\subseteq H_{L'}\bmod K$ and $m(L)\ge m(L')$. A twin catch
at phase $\varphi$ for the truncated twin implies a catch for the full twin, so
$\{M=1,M'_L=0\}\subseteq\{M=1,M'=0\}$ and $D_L\le D$; Lemma~\ref{lem:c3} holds for any pairing.
\end{proof}

\begin{remark}[check]
On the twin of the 78 static scenes, shortening every event by $x\in\{10,25,50\}\%$ gives non-decreasing $q$ for every
event and every $K\in\{1,4,8,16,48\}$ (all monotone: \numCfiveAll); e.g.\ at $d_{\min}=32$, $K=16$, $q$ rises from
\numCfiveQzero{} to \numCfiveQfifty{} at $x=50\%$.
\end{remark}


}

% ------------------------------------------------------------------------------------------------------
\section{Units and event counts}\label{sec:units}
\paragraph{Unit.} Every rate is per \emph{event} (a maximal run of frames containing the ground-truth object whose
largest instance has at least 256 pixels; events shorter than $d_{\min}$ frames are dropped). Each event has three twin
renderings (sprite variants); the twin catches the event at a sampling phase if at least two variants catch it. Under a
uniform phase the discordance probability of an event is $x_i=|R'_i\setminus R_i|/K$; the certificates use the
integer worst-phase count $W_i=\mathbf 1\{R'_i\setminus R_i\ne\emptyset\}\ge x_i$ (remark on the event unit in Section~\ref{sec:proofs}), and the phase-averaged count is reported in parentheses. Table~\ref{tab:rules} shows that the variant
rule (1, 2 or 3 of 3) does not change any static-domain count. Table~\ref{tab:waterfall} reconciles the event counts
of the initial twin run with the analysed set.


% ------------------------------------------------------------------------------------------------------
\section{Full operating-point grid}
Table~\ref{tab:grid} lists all \nGridPts{} operating points at $\varepsilon=20\%$. The twin route certifies a camera with
no real event when $q_{\mathrm{U}}+\UCP_{\mathrm{pop}}\le\varepsilon$ with $\delta$ split in half, $\UCP_{\mathrm{pop}}$
computed without the camera. ``As built'' uses the twin's own miss rate bounded from three runs per event; ``converged''
treats $q$ as known. Direct certification needs $n_{\mathrm{dir}}$ events at 80\% power.

\section{Audit details}
Table~\ref{tab:loso} lists the scenes whose discordance exceeds the leave-one-scene-out bound. Table~\ref{tab:mf1} gives
must-fail scenario 1 (temporal) and Tables~\ref{tab:nightOP} and~\ref{tab:leakOP} must-fail scenario 2 (night scenes
leaked into the calibration set), including the per-camera out-of-domain pass-through; Table~\ref{tab:ood} judges
the camera-jitter and night cameras with a bound from the static scenes under the worst-phase convention. The static
domain is defined by the category labels of the datasets (CDnet without \emph{PTZ}, \emph{cameraJitter},
\emph{nightVideos} and \emph{turbulence}; LASIESTA without SM and MC); this definition was fixed after the initial
version of this study, and Table~\ref{tab:v10v11} compares the static domain of the three versions. The camera-jitter
tier holds the CDnet \emph{cameraJitter} scenes and the LASIESTA scenes with simulated camera motion (SM); its only
discordant events and its only false acceptance are in \emph{traffic}, a CDnet scene with real shake. No static scene has 59 or more
qualifying events, so the ``truth'' of a camera is its own finite event pool (exact over phases).

\section{Sensitivity}
Table~\ref{tab:iw} compares the synthetic-versus-real classifier on all target sets (main text: \nIWtemporal,
\nIWgeometry, \nIWappearance) and on CDnet targets only (\nIWtemporalCD, \nIWgeometryCD, \nIWappearanceCD).
Table~\ref{tab:aspect} gives twin and real hit rates on the same frames by sprite aspect error; the within-event paired
difference on clamped frames is \nClampExtraDrop{} points over \nClampPairs{} event--variant pairs.
Table~\ref{tab:c5a} checks C5a and Table~\ref{tab:c2emp} the empirical C2 procedure with every simulator.

\section{Twin construction and licences}\label{sec:twin}
\paragraph{Construction.} Per scene, the background plate is the per-pixel median of up to 60 frames without ground-truth
foreground (longest empty run preferred). Per event and variant (seeds 42, 43, 44): the real object (ground-truth mask,
dilated 3 px) is replaced by the plate; the sprite height equals the ground-truth box height, its width is the sprite
aspect times the height clamped to $\pm10\%$ of the box width, bottom-centred on the box; brightness gain is the frame
luminance over a reference (clipped to $[0.35,1.2]$) with a 5\% blend of the frame mean colour; a soft elliptical shadow
and a 2-px feather are applied. Person sprites follow a four-frame walking cycle, mirrored to the motion direction;
vehicle sprites are chosen by aspect ratio. The detector is run on every twin frame until the hit residues cover all 48
classes modulo 48 (exact early stop for every $K$ dividing 48).
\paragraph{Licences.} Vehicle sprites were cut from Virtual KITTI~2 (Cabon et al., 2020; CC BY-NC-SA 3.0) and are not
redistributed; the extraction script is provided. Person sprites are generated with MPFB~2 from its CC0 bundled assets.
No frame of CDnet~2014 or LASIESTA is reproduced in the paper or this supplement, because no reuse terms for figures
could be recorded; the illustration of sprites uses CC0 assets on a neutral background.

\section{Supplementary tables}
% generated by tools/make_supp_tables.py from results/*.json -- DO NOT EDIT
\begin{table}[H]
\centering
\caption{Event counts from the initial twin run to the analysed set (primary definition: largest object $\ge 256$ px). The two headline counts differ in unit: all durations vs.\ duration $\ge 32$.}
\label{tab:waterfall}
\small
\begin{tabular}{@{}lrr@{}}
\toprule
step & events & scenes \\
\midrule
initial twin run, original definition (all events, all durations) & 329 & 64 \\
initial twin run, largest object $\ge 256$ px & 257 & 64 \\
same scenes, LASIESTA frames in true temporal order & 257 & 64 \\
LASIESTA scene list rebuilt (20 to 48 scenes) & 308 & 92 \\
pan-tilt-zoom scenes excluded (4) & 283 & 88 \\
LASIESTA moving-camera scenes excluded (4) & 279 & 84 \\
scenes run (2 held-out turbulence scenes not run) & 269 & 82 \\
events of at least $d_{\min}=32$ frames & 203 & 82 \\
\bottomrule
\end{tabular}
\end{table}

\begin{table}[H]
\centering
\caption{Sprite-variant rule (twin catches if $\ge r$ of 3 variants catch) at OP$^\ast$: worst-phase discordant events; phase-averaged count under the rule (exp.) and averaged over the three variants (var.\ mean, the parenthesised count of the main text); bounds.}
\label{tab:rules}
\small
\begin{tabular}{@{}llrrrrrrr@{}}
\toprule
tier & $r$ & $N$ & disc. & exp. & var.\ mean & $U(X,N)$ & $J/m$ & $U(J,m)$ \\
\midrule
main & 1/3 & 114 & 0 & 0.00 & 0.00 & 2.59\% & 0/48 & 6.05\% \\
main & 2/3 & 114 & 0 & 0.00 & 0.00 & 2.59\% & 0/48 & 6.05\% \\
main & 3/3 & 114 & 0 & 0.00 & 0.00 & 2.59\% & 0/48 & 6.05\% \\
jitter & 1/3 & 62 & 3 & 0.31 & 0.31 & 12.03\% & 1/26 & 16.98\% \\
jitter & 2/3 & 62 & 3 & 0.31 & 0.31 & 12.03\% & 1/26 & 16.98\% \\
jitter & 3/3 & 62 & 3 & 0.31 & 0.31 & 12.03\% & 1/26 & 16.98\% \\
night & 1/3 & 25 & 13 & 8.75 & 8.21 & 69.49\% & 5/6 & 99.15\% \\
night & 2/3 & 25 & 12 & 8.31 & 8.21 & 65.86\% & 5/6 & 99.15\% \\
night & 3/3 & 25 & 12 & 7.56 & 8.21 & 65.86\% & 5/6 & 99.15\% \\
turbulence & 1/3 & 2 & 1 & 0.12 & 0.12 & 97.47\% & 1/2 & 97.47\% \\
turbulence & 2/3 & 2 & 1 & 0.12 & 0.12 & 97.47\% & 1/2 & 97.47\% \\
turbulence & 3/3 & 2 & 1 & 0.12 & 0.12 & 97.47\% & 1/2 & 97.47\% \\
all tiers & 1/3 & 203 & 17 & 9.19 & 8.65 & 12.30\% & 7/82 & 15.44\% \\
all tiers & 2/3 & 203 & 16 & 8.75 & 8.65 & 11.73\% & 7/82 & 15.44\% \\
all tiers & 3/3 & 203 & 16 & 8.00 & 8.65 & 11.73\% & 7/82 & 15.44\% \\
\bottomrule
\end{tabular}
\end{table}

\begingroup\footnotesize
\begin{longtable}{@{}rrrrrrrr@{}}
\caption{Full operating-point grid at $\varepsilon=20\%$ (48 static cameras): static-domain real miss rate, events for direct certification, median events saved per camera (twin run to convergence), cameras saving (converged / as built), CPU-hours per saved event.}\label{tab:grid}\\
\toprule
$d_{\min}$ & $K$ & $\hat p$ & $n_{\mathrm{dir}}$ & saved & cams & as built & CPU-h \\
\midrule
\endfirsthead
\toprule
$d_{\min}$ & $K$ & $\hat p$ & $n_{\mathrm{dir}}$ & saved & cams & as built & CPU-h \\
\midrule
\endhead
\bottomrule
\endlastfoot
32 & 1 & 0.88\% & 14 & 14 & 47 & 3 & 0.064 \\
32 & 2 & 0.88\% & 14 & 14 & 47 & 3 & 0.064 \\
32 & 3 & 0.88\% & 14 & 14 & 47 & 3 & 0.064 \\
32 & 4 & 0.88\% & 14 & 14 & 47 & 3 & 0.064 \\
32 & 6 & 0.88\% & 14 & 14 & 47 & 3 & 0.064 \\
32 & 8 & 0.88\% & 14 & 14 & 47 & 3 & 0.064 \\
32 & 12 & 0.88\% & 14 & 14 & 47 & 3 & 0.070 \\
32 & 16 & 0.88\% & 14 & 14 & 46 & 3 & 0.069 \\
32 & 24 & 1.21\% & 14 & 14 & 44 & 0 & 0.351 \\
32 & 48 & 5.94\% & 37 & 0 & 0 & 0 & -- \\
24 & 1 & 0.84\% & 14 & 14 & 47 & 3 & 0.064 \\
24 & 2 & 0.84\% & 14 & 14 & 47 & 3 & 0.064 \\
24 & 3 & 0.84\% & 14 & 14 & 47 & 3 & 0.064 \\
24 & 4 & 0.84\% & 14 & 14 & 47 & 3 & 0.064 \\
24 & 6 & 0.84\% & 14 & 14 & 47 & 3 & 0.064 \\
24 & 8 & 1.05\% & 14 & 14 & 47 & 0 & 0.094 \\
24 & 12 & 1.12\% & 14 & 14 & 45 & 0 & 0.160 \\
24 & 16 & 1.68\% & 22 & 22 & 45 & 0 & 0.106 \\
24 & 24 & 2.28\% & 22 & 22 & 43 & 0 & 0.229 \\
24 & 48 & 8.23\% & 57 & 0 & 0 & 0 & -- \\
16 & 1 & 1.55\% & 14 & 14 & 47 & 4 & 0.064 \\
16 & 2 & 1.94\% & 22 & 22 & 47 & 3 & 0.041 \\
16 & 3 & 2.07\% & 22 & 22 & 47 & 3 & 0.041 \\
16 & 4 & 2.13\% & 22 & 22 & 47 & 3 & 0.043 \\
16 & 6 & 2.20\% & 22 & 22 & 46 & 3 & 0.042 \\
16 & 8 & 2.62\% & 22 & 22 & 46 & 0 & 0.061 \\
16 & 12 & 3.42\% & 22 & 22 & 45 & 0 & 0.102 \\
16 & 16 & 4.26\% & 30 & 30 & 45 & 0 & 0.078 \\
16 & 24 & 6.07\% & 37 & 37 & 43 & 0 & 0.136 \\
16 & 48 & 13.45\% & 210 & 0 & 0 & 0 & -- \\
12 & 1 & 2.31\% & 22 & 22 & 47 & 4 & 0.041 \\
12 & 2 & 2.69\% & 22 & 22 & 47 & 3 & 0.041 \\
12 & 3 & 2.82\% & 22 & 22 & 47 & 3 & 0.041 \\
12 & 4 & 2.88\% & 22 & 22 & 47 & 3 & 0.043 \\
12 & 6 & 2.95\% & 22 & 22 & 46 & 3 & 0.042 \\
12 & 8 & 3.37\% & 22 & 22 & 46 & 0 & 0.061 \\
12 & 12 & 4.17\% & 30 & 30 & 45 & 0 & 0.074 \\
12 & 16 & 5.00\% & 30 & 30 & 45 & 0 & 0.078 \\
12 & 24 & 6.79\% & 44 & 44 & 43 & 0 & 0.115 \\
12 & 48 & 14.12\% & 260 & 0 & 0 & 0 & -- \\
8 & 1 & 2.22\% & 22 & 22 & 47 & 4 & 0.041 \\
8 & 2 & 2.59\% & 22 & 22 & 47 & 4 & 0.041 \\
8 & 3 & 2.72\% & 22 & 22 & 47 & 3 & 0.043 \\
8 & 4 & 2.96\% & 22 & 22 & 47 & 0 & 0.063 \\
8 & 6 & 2.96\% & 22 & 22 & 46 & 3 & 0.042 \\
8 & 8 & 3.98\% & 30 & 30 & 46 & 0 & 0.073 \\
8 & 12 & 5.56\% & 37 & 37 & 44 & 0 & 0.062 \\
8 & 16 & 6.90\% & 44 & 44 & 44 & 0 & 0.053 \\
8 & 24 & 9.17\% & 69 & 69 & 43 & 0 & 0.073 \\
8 & 48 & 16.76\% & 901 & 0 & 0 & 0 & -- \\
4 & 1 & 2.72\% & 22 & 22 & 47 & 4 & 0.041 \\
4 & 2 & 3.40\% & 22 & 22 & 47 & 3 & 0.043 \\
4 & 3 & 4.08\% & 30 & 30 & 46 & 3 & 0.031 \\
4 & 4 & 4.76\% & 30 & 30 & 46 & 0 & 0.045 \\
4 & 6 & 6.35\% & 44 & 44 & 46 & 0 & 0.031 \\
4 & 8 & 8.42\% & 63 & 63 & 45 & 0 & 0.035 \\
4 & 12 & 11.00\% & 106 & 106 & 44 & 0 & 0.021 \\
4 & 16 & 12.80\% & 170 & 170 & 44 & 0 & 0.014 \\
4 & 24 & 15.45\% & 441 & 441 & 43 & 0 & 0.011 \\
4 & 48 & 22.99\% & -- & 0 & 0 & 0 & -- \\
1 & 1 & 3.11\% & 22 & 22 & 47 & 0 & 0.063 \\
1 & 2 & 5.90\% & 37 & 37 & 46 & 0 & 0.059 \\
1 & 3 & 7.66\% & 50 & 50 & 46 & 0 & 0.044 \\
1 & 4 & 9.47\% & 76 & 76 & 45 & 0 & 0.029 \\
1 & 6 & 12.11\% & 135 & 135 & 45 & 0 & 0.017 \\
1 & 8 & 14.60\% & 310 & 310 & 45 & 0 & 0.007 \\
1 & 12 & 17.55\% & 1583 & 1583 & 44 & 0 & 0.001 \\
1 & 16 & 19.49\% & -- & 0 & 0 & 0 & -- \\
1 & 24 & 22.20\% & -- & 0 & 0 & 0 & -- \\
1 & 48 & 29.39\% & -- & 0 & 0 & 0 & -- \\
\end{longtable}
\endgroup
\begin{table}[H]
\centering
\caption{Night scenes at OP$^\ast$: real miss rate split into detector and temporal misses, twin miss rate, discordance, and whether the camera is falsely certified when judged with a bound from the static scenes only.}
\label{tab:nightOP}
\small
\begin{tabular}{@{}lrrrrrrrr@{}}
\toprule
scene & $n$ & $p_s$ & detector & temporal & $q_s$ & $D_s$ & disc. & false cert. \\
\midrule
bridgeEntry & 2 & 15.6\% & 15.6\% & 0.0\% & 0.0\% & 15.6\% & 1 & no \\
busyBoulvard & 4 & 96.9\% & 96.9\% & 0.0\% & 0.0\% & 96.9\% & 4 & yes \\
fluidHighway & 4 & 26.6\% & 26.6\% & 0.0\% & 0.0\% & 26.6\% & 2 & yes \\
streetCornerAtNight & 10 & 28.1\% & 28.1\% & 0.0\% & 11.9\% & 20.6\% & 4 & yes \\
tramStation & 2 & 0.0\% & 0.0\% & 0.0\% & 0.0\% & 0.0\% & 0 & no \\
winterStreet & 3 & 33.3\% & 33.3\% & 0.0\% & 0.0\% & 33.3\% & 1 & yes \\
\bottomrule
\end{tabular}
\end{table}

\begin{table}[H]
\centering
\caption{Must-fail 2 at OP$^\ast$: $k$ night scenes in the calibration set (all subsets).}
\label{tab:leakOP}
\small
\begin{tabular}{@{}rrrrrrr@{}}
\toprule
$k$ & subsets & $U_{\mathrm{pop}}$ med. & range & $U_{\mathrm{fleet}}$ med. & withdrawn & static false cert. \\
\midrule
0 & 1 & 6.1\% & 6.1\%--6.1\% & 2.6\% & 0\% & 0 \\
1 & 6 & 9.3\% & 5.9\%--9.3\% & 4.6\% & 0\% & 0 \\
2 & 15 & 12.1\% & 9.1\%--12.1\% & 7.5\% & 67\% & 0 \\
3 & 20 & 13.2\% & 11.8\%--14.5\% & 9.1\% & 100\% & 0 \\
4 & 15 & 14.2\% & 14.2\%--16.7\% & 10.6\% & 100\% & 0 \\
5 & 6 & 16.4\% & 16.4\%--18.8\% & 12.6\% & 100\% & 0 \\
6 & 1 & 18.5\% & 18.5\%--18.5\% & 13.6\% & 100\% & 0 \\
\bottomrule
\end{tabular}
\end{table}

\begin{table}[H]
\centering
\caption{Must-fail 1 (temporal, $K=48$): static scenes with $p_s>\varepsilon$; detector vs.\ temporal misses and twin miss rate. All were refused in every Monte-Carlo run.}
\label{tab:mf1}
\small
\begin{tabular}{@{}llrrrrr@{}}
\toprule
point & scene & $n$ & $p_s$ & detector & temporal & $q_s$ \\
\midrule
d1 K48 & port\_0\_17fps & 5 & 99.6\% & 43.3\% & 56.2\% & 100.0\% \\
d1 K48 & tramCrossroad\_1fps & 11 & 77.7\% & 6.1\% & 71.6\% & 74.2\% \\
d1 K48 & tunnelExit\_0\_35fps & 34 & 83.6\% & 5.2\% & 78.4\% & 87.9\% \\
d4 K48 & port\_0\_17fps & 5 & 99.6\% & 43.3\% & 56.2\% & 100.0\% \\
d4 K48 & tramCrossroad\_1fps & 8 & 71.1\% & 8.1\% & 63.0\% & 66.7\% \\
d4 K48 & tunnelExit\_0\_35fps & 23 & 77.3\% & 7.5\% & 69.7\% & 83.4\% \\
d16 K48 & port\_0\_17fps & 3 & 99.3\% & 57.6\% & 41.7\% & 100.0\% \\
d16 K48 & tramCrossroad\_1fps & 4 & 55.7\% & 12.5\% & 43.2\% & 47.9\% \\
d16 K48 & tunnelExit\_0\_35fps & 11 & 61.7\% & 11.9\% & 49.8\% & 72.7\% \\
\bottomrule
\end{tabular}
\end{table}

\begin{table}[H]
\centering
\caption{Leave-one-scene-out: scenes whose discordance exceeds the bound computed without them.}
\label{tab:loso}
\small
\begin{tabular}{@{}llrrrrrrr@{}}
\toprule
point & scene & $n$ & $D_s$ & $p_s$ & $q_s$ & $U_{-s}$ & $p_s>q_s+U$ & false cert. \\
\midrule
-- & none &  &  &  &  &  &  &  \\
\bottomrule
\end{tabular}
\end{table}

\begin{table}[H]
\centering
\caption{Synthetic-vs-real classifier (logistic), median AUC over target sets: all targets (main text) and CDnet targets only.}
\label{tab:iw}
\small
\begin{tabular}{@{}lrr@{}}
\toprule
features & all targets & CDnet targets \\
\midrule
temporal & 0.64 & 0.61 \\
geometry & 0.94 & 0.96 \\
appearance only & 0.77 & 0.75 \\
\bottomrule
\end{tabular}
\end{table}

\begin{table}[H]
\centering
\caption{Twin and real hit rates on the same frames by object kind and sprite aspect error (before the width clamp).}
\label{tab:aspect}
\small
\begin{tabular}{@{}llrrrr@{}}
\toprule
kind & $|$aspect err.$|$ & frames & twin hit & real hit & diff. (pts) \\
\midrule
person & $\le$10\% & 8870 & 91.7\% & 94.8\% & -3.2 \\
person & 10-20\% & 5899 & 90.1\% & 94.2\% & -4.2 \\
person & 20-40\% & 4404 & 74.0\% & 86.2\% & -12.2 \\
person & $>$40\% & 5449 & 34.4\% & 44.5\% & -10.0 \\
vehicle & $\le$10\% & 4105 & 86.9\% & 89.0\% & -2.1 \\
vehicle & 10-20\% & 2569 & 88.0\% & 90.2\% & -2.2 \\
vehicle & 20-40\% & 3087 & 87.6\% & 85.1\% & 2.5 \\
vehicle & $>$40\% & 3881 & 64.9\% & 74.2\% & -9.3 \\
\bottomrule
\end{tabular}
\end{table}

\begin{table}[H]
\centering
\caption{C5a: twin miss rate when every event is shortened by $x$.}
\label{tab:c5a}
\small
\begin{tabular}{@{}rrrrrrr@{}}
\toprule
$d_{\min}$ & $K$ & $x=0$ & 10\% & 25\% & 50\% & per event monotone \\
\midrule
32 & 1 & 0.88\% & 0.88\% & 0.88\% & 0.88\% & yes \\
32 & 8 & 0.99\% & 0.99\% & 0.99\% & 1.21\% & yes \\
32 & 16 & 1.37\% & 1.43\% & 1.59\% & 2.36\% & yes \\
32 & 48 & 7.05\% & 7.80\% & 10.29\% & 19.85\% & yes \\
16 & 1 & 2.33\% & 2.33\% & 2.33\% & 2.33\% & yes \\
16 & 8 & 4.36\% & 4.46\% & 4.65\% & 5.72\% & yes \\
16 & 16 & 6.54\% & 6.73\% & 7.17\% & 9.16\% & yes \\
16 & 48 & 15.04\% & 15.94\% & 18.56\% & 27.65\% & yes \\
1 & 1 & 4.97\% & 4.97\% & 5.59\% & 6.83\% & yes \\
1 & 8 & 16.69\% & 16.85\% & 18.01\% & 21.04\% & yes \\
1 & 16 & 21.66\% & 21.89\% & 22.83\% & 25.50\% & yes \\
1 & 48 & 30.77\% & 31.52\% & 33.81\% & 41.46\% & yes \\
\bottomrule
\end{tabular}
\end{table}

\begingroup\footnotesize
\begin{table}[H]
\centering
\caption{Empirical C2: two-tier procedure with each simulator (100{,}000 runs), $\delta'=7.5/10/15/20\%$.}
\label{tab:c2emp}
\begin{tabular}{@{}llrrll@{}}
\toprule
point & simulator & $N_{\mathrm{sim}}$ & pass & credit (events) & false cert.\ at $p=\varepsilon$ \\
\midrule
$\varepsilon=5\%$, (32,1) & BMC (G$^\ast$) & 30 & 0.0\% & 0.0/0.0/0.0/0.0 & 4.9\%/4.8\%/4.8\%/4.8\% \\
$\varepsilon=5\%$, (32,1) & BMC (temporal) & 49 & 0.0\% & 0.0/0.0/0.0/0.0 & 4.8\%/4.8\%/4.7\%/4.9\% \\
$\varepsilon=5\%$, (32,1) & paired twin & 114 & 73.3\% & 5.9/10.3/16.1/19.8 & 6.7\%/8.4\%/12.2\%/15.6\% \\
$\varepsilon=20\%$, OP$^\ast$ & BMC (G$^\ast$) & 30 & 100.0\% & 2.0/3.0/5.0/6.0 & 7.0\%/8.5\%/13.4\%/16.7\% \\
$\varepsilon=20\%$, OP$^\ast$ & BMC (temporal) & 49 & 100.0\% & 2.0/3.0/5.0/6.0 & 6.9\%/8.6\%/13.4\%/17.0\% \\
$\varepsilon=20\%$, OP$^\ast$ & paired twin & 114 & 100.0\% & 2.0/3.0/5.0/6.0 & 6.8\%/8.5\%/13.5\%/16.8\% \\
\bottomrule
\end{tabular}
\end{table}
\endgroup
\begingroup\footnotesize\setlength{\tabcolsep}{2pt}
\begin{table}[H]
\centering
\caption{Worst-phase convention (valid for any fixed deployment phase): real and twin miss rates averaged over phases and at the worst phase, discordant events, fleet and population bounds.}
\label{tab:worstphase}
\begin{tabular}{@{}llrrrrrrrrr@{}}
\toprule
point & tier & $N$ & $\hat p$ avg & $\hat p$ worst & $\hat q$ avg & $\hat q$ worst & disc. & $U(X,N)$ & $J/m$ & $U(J,m)$ \\
\midrule
OP$^\ast$ & main & 114 & 0.9\% & 0.9\% & 1.4\% & 3.5\% & 0 & 2.59\% & 0/48 & 6.05\% \\
 & jitter & 62 & 0.5\% & 4.8\% & 0.0\% & 0.0\% & 3 & 12.03\% & 1/26 & 16.98\% \\
 & night & 25 & 36.2\% & 52.0\% & 4.8\% & 16.0\% & 12 & 65.86\% & 5/6 & 99.15\% \\
 & turbulence & 2 & 6.2\% & 50.0\% & 0.0\% & 0.0\% & 1 & 97.47\% & 1/2 & 97.47\% \\
 & all tiers & 203 & 5.2\% & 8.9\% & 1.4\% & 3.9\% & 16 & 11.73\% & 7/82 & 15.44\% \\
\bottomrule
\end{tabular}
\end{table}
\endgroup
\begin{table}[H]
\centering
\caption{Exchange rate under the worst-phase convention (48 static cameras; twin run to convergence).}
\label{tab:worstx}
\small
\begin{tabular}{@{}lrrrrrr@{}}
\toprule
point & $\hat p$ worst & $n_{\mathrm{dir}}$ & accepted & saved & false acc. & twin runs \\
\midrule
OP$^\ast$ & 0.9\% & 14 & 44 & 14 & 0 & 28 \\
\bottomrule
\end{tabular}
\end{table}

\begingroup\footnotesize\setlength{\tabcolsep}{2pt}
\begin{table}[H]
\centering
\caption{Out-of-domain cameras (camera jitter, night) judged by the twin route calibrated on the static scenes only (worst-phase convention, transfer term $U_{\delta/2}(J,m)$ of the static scenes): accepted, and falsely accepted ($p_w>\varepsilon$).}
\label{tab:ood}
\begin{tabular}{@{}lllrrrrll@{}}
\toprule
point & tier & camera & $n$ & $p_w$ & $q_w$ & $U_{\delta/2}$ & accepted & false acc. \\
\midrule
OP$^\ast$ & jitter & I\_SM\_01 & 3 & 0.0\% & 0.0\% & 7.4\% & yes & no \\
 & jitter & I\_SM\_02 & 2 & 0.0\% & 0.0\% & 7.4\% & yes & no \\
 & jitter & I\_SM\_03 & 2 & 0.0\% & 0.0\% & 7.4\% & yes & no \\
 & jitter & I\_SM\_04 & 3 & 0.0\% & 0.0\% & 7.4\% & yes & no \\
 & jitter & I\_SM\_05 & 3 & 0.0\% & 0.0\% & 7.4\% & yes & no \\
 & jitter & I\_SM\_06 & 3 & 0.0\% & 0.0\% & 7.4\% & yes & no \\
 & jitter & I\_SM\_07 & 3 & 0.0\% & 0.0\% & 7.4\% & yes & no \\
 & jitter & I\_SM\_08 & 3 & 0.0\% & 0.0\% & 7.4\% & yes & no \\
 & jitter & I\_SM\_09 & 3 & 0.0\% & 0.0\% & 7.4\% & yes & no \\
 & jitter & I\_SM\_10 & 3 & 0.0\% & 0.0\% & 7.4\% & yes & no \\
 & jitter & I\_SM\_11 & 3 & 0.0\% & 0.0\% & 7.4\% & yes & no \\
 & jitter & I\_SM\_12 & 3 & 0.0\% & 0.0\% & 7.4\% & yes & no \\
 & jitter & O\_SM\_01 & 1 & 0.0\% & 0.0\% & 7.4\% & yes & no \\
 & jitter & O\_SM\_02 & 1 & 0.0\% & 0.0\% & 7.4\% & yes & no \\
 & jitter & O\_SM\_03 & 1 & 0.0\% & 0.0\% & 7.4\% & yes & no \\
 & jitter & O\_SM\_04 & 1 & 0.0\% & 0.0\% & 7.4\% & yes & no \\
 & jitter & O\_SM\_05 & 1 & 0.0\% & 0.0\% & 7.4\% & yes & no \\
 & jitter & O\_SM\_06 & 1 & 0.0\% & 0.0\% & 7.4\% & yes & no \\
 & jitter & O\_SM\_07 & 1 & 0.0\% & 0.0\% & 7.4\% & yes & no \\
 & jitter & O\_SM\_08 & 1 & 0.0\% & 0.0\% & 7.4\% & yes & no \\
 & jitter & O\_SM\_09 & 1 & 0.0\% & 0.0\% & 7.4\% & yes & no \\
 & jitter & O\_SM\_10 & 1 & 0.0\% & 0.0\% & 7.4\% & yes & no \\
 & jitter & O\_SM\_11 & 1 & 0.0\% & 0.0\% & 7.4\% & yes & no \\
 & jitter & O\_SM\_12 & 1 & 0.0\% & 0.0\% & 7.4\% & yes & no \\
 & jitter & boulevard & 8 & 0.0\% & 0.0\% & 7.4\% & yes & no \\
 & jitter & traffic & 8 & 37.5\% & 0.0\% & 7.4\% & yes & yes \\
 & night & bridgeEntry & 2 & 50.0\% & 0.0\% & 7.4\% & yes & yes \\
 & night & busyBoulvard & 4 & 100.0\% & 0.0\% & 7.4\% & yes & yes \\
 & night & fluidHighway & 4 & 50.0\% & 0.0\% & 7.4\% & yes & yes \\
 & night & streetCornerAtNight & 10 & 50.0\% & 40.0\% & 7.4\% & no & no \\
 & night & tramStation & 2 & 0.0\% & 0.0\% & 7.4\% & yes & no \\
 & night & winterStreet & 3 & 33.3\% & 0.0\% & 7.4\% & yes & yes \\
\bottomrule
\end{tabular}
\end{table}
\endgroup
\begin{table}[H]
\centering
\caption{Site-level population bound at OP$^\ast$: 17 sites, $J_{\mathrm{site}}=0$, $U(J_{\mathrm{site}},m_{\mathrm{site}})=16.2\%$; per site: scenes, events, discordant events, bound without the site, leave-one-site-out violation.}
\label{tab:siteOP}
\small
\begin{tabular}{@{}lrrrrl@{}}
\toprule
site & scenes & events & disc. & $U_{-\mathrm{site}}$ & violation \\
\midrule
CDnet badWeather & 4 & 15 & 0 & 17.1\% & no \\
CDnet baseline & 1 & 2 & 0 & 17.1\% & no \\
CDnet dynamicBackground & 6 & 19 & 0 & 17.1\% & no \\
CDnet intermittentObjectMotion & 3 & 4 & 0 & 17.1\% & no \\
CDnet lowFramerate & 4 & 5 & 0 & 17.1\% & no \\
CDnet shadow & 6 & 37 & 0 & 17.1\% & no \\
CDnet thermal & 4 & 7 & 0 & 17.1\% & no \\
LASIESTA I\_BS & 2 & 2 & 0 & 17.1\% & no \\
LASIESTA I\_CA & 2 & 2 & 0 & 17.1\% & no \\
LASIESTA I\_IL & 2 & 2 & 0 & 17.1\% & no \\
LASIESTA I\_MB & 2 & 2 & 0 & 17.1\% & no \\
LASIESTA I\_OC & 2 & 4 & 0 & 17.1\% & no \\
LASIESTA I\_SI & 2 & 4 & 0 & 17.1\% & no \\
LASIESTA O\_CL & 2 & 2 & 0 & 17.1\% & no \\
LASIESTA O\_RA & 2 & 3 & 0 & 17.1\% & no \\
LASIESTA O\_SN & 2 & 2 & 0 & 17.1\% & no \\
LASIESTA O\_SU & 2 & 2 & 0 & 17.1\% & no \\
\bottomrule
\end{tabular}
\end{table}

\begin{table}[H]
\centering
\caption{Static domain at OP$^\ast$ in the three versions of this study: 78 scenes; 76 scenes (CDnet camera-jitter scenes moved to the jitter tier); 48 scenes (LASIESTA simulated-motion scenes moved to the jitter tier and LASIESTA moving-camera scenes excluded; this version). Each column is recomputed by the released code from the same twin run and, for the two earlier versions, equals the published numbers; the last column gives the scenes that carry the earlier disagreements.}
\label{tab:v10v11}
\small
\begin{tabular}{@{}lllll@{}}
\toprule
quantity & 78 scenes & 76 scenes & 48 scenes & located in \\
\midrule
static scenes / events & 78 / 180 & 76 / 164 & 48 / 114 & -- \\
discordant events (worst phase), $U(X,N)$ & 3, 4.25\% & 0, 1.81\% & 0, 2.59\% & traffic \\
discordant scenes $J/m$, $U(J,m)$ & 1/78, 5.94\% & 0/76, 3.87\% & 0/48, 6.05\% & traffic \\
scenes above their LOSO bound & 1/78 & 0/76 & 0/48 & traffic \\
worst-phase exchange: $n_{\mathrm{dir}}$; accepted; false acc. & 22; 74/78; 1 & 14; 72/76; 0 & 14; 44/48; 0 & traffic \\
\bottomrule
\end{tabular}
\end{table}


\bibliographystyle{elsarticle-num}
\bibliography{refs_ress}

\end{document}
